% Generated draft: requires mathematical and source review.
\documentclass[11pt]{article}
\usepackage[margin=25mm]{geometry}
\usepackage{fontspec}
\setmainfont{Latin Modern Roman}
\usepackage{amsmath,amssymb,mathtools,unicode-math}
\setmathfont{Latin Modern Math}
\usepackage{enumitem,needspace,xurl,hyperref}
\setlength{\parindent}{0pt}
\setlength{\parskip}{5pt}
\setlength{\emergencystretch}{3em}
\newcommand{\N}{\mathbb{N}}
\newcommand{\Z}{\mathbb{Z}}
\newcommand{\Q}{\mathbb{Q}}
\newcommand{\R}{\mathbb{R}}
\newcommand{\C}{\mathbb{C}}
\newcommand{\F}{\mathbb{F}}
\newcommand{\A}{\mathbb{A}}
\newcommand{\PP}{\mathbb{P}}
\newcommand{\E}{\mathbb{E}}
\newcommand{\eps}{\varepsilon}
\newcommand{\dd}{\,\mathrm{d}}
\begin{document}
\section*{UnsolvedMath ID: 11000318}
\textbf{Source problem code:} catalogue entry
\subsection{Definitions and mathematical statement}
Retain the source Lie algebra \(\mathfrak h_{g,1}^{\Q}\), its symplectic-group action \(\operatorname{Sp}\), and the classes
\[
 e_1,t_3,t_5,\ldots\ \in\
 V_g=H^2(\mathfrak h_{g,1}^{\Q};\Q)^{\operatorname{Sp}}.
\]
Here the superscript denotes invariant classes in Lie-algebra cohomology. The precise Lie algebra, coefficient convention, cochain/support convention and class constructions occur in missing preceding text; the usual symplectic-derivation interpretation is not imposed as a replacement definition.

The conjecture asserts all of these classes are nonzero, linearly independent and span \(V_g\). In an ordinary vector-space basis interpretation, every relation
\[
 a_0e_1+\sum_{j=1}^{r}a_jt_{2j+1}=0,
 \qquad a_j\in\Q,
\]
has all coefficients zero, and every class in \(V_g\) is a finite such linear combination. If a completed/continuous cohomology or a weightwise stable interpretation is intended, the corresponding meaning of basis must instead be specified. The source does not give a genus range or say whether the claim is for fixed \(g\) or the stable limit. No formula for the unnamed classes or partial nonvanishing result is invented.
\subsection{Short English statement}
The source conjectures that \(e_1,t_3,t_5,\ldots\) are nonzero and form a basis of the symplectic-invariant second cohomology of the indicated Lie algebra, with genus and completion conventions unresolved.
\subsection{Sources}
\begin{itemize}
\item \emph{Problems on Mapping Class Groups and Related Topics}, edited by Farb, Conjecture 6.1, supplied PDF page 371. \url{https://www.math.uchicago.edu/~farb/papers/mcgbook.pdf}. Recorded access 2026-07-29; extraction pdf-font-structure. The preceding Lie-algebra and class definitions are absent.
\item Local UnsolvedMath registry, \texttt{data/unsolvedmath/problems.json}, permanent ID \(11000318\), exact source code \texttt{catalogue entry}; supplied OPEN-TRIAGE summary mentions partial Morita-type nontriviality without a precise theorem or bibliography.
\end{itemize}
\end{document}
